\vspace{-1em}
\section{Introduction}

Recent advances in talking head generation have greatly improved lip synchronization, identity preservation, and photorealistic rendering~\cite{grassal2022neural,ma2024cvthead,prajwal2020lip,ji2021audio,liang2022expressive,zhang2023sadtalker,geng2023realtalk,zhu2025infp}. Most existing methods, however, focus on animating the speaker. In face-to-face conversation, the listener is equally important: non-verbal behaviors, such as brief reactions like smiling and nodding, convey attention, agreement, hesitation, and affective engagement~\cite{wang2025diffusion}. The timing and type of these reactions strongly affect perceived naturalness, since visually plausible motions can still appear unnatural when reactions are absent, mistimed, or inconsistent with the conversational context~\cite{murray2022learning,ng2022learning,de2012survey}.

Listening head generation has recently emerged as a distinct task. ViCo~\cite{zhou2022vico} introduced an early benchmark, while subsequent works explored non-deterministic listener motion, language-aware responses, emotion conditioning, improved sequence modeling, and real-time dyadic interaction~\cite{ng2022learning,ng2023can,song2023emotional,liu2024listenformer,zhu2025infp,guo2025arig}. Despite this progress, reaction-aware listening head generation remains limited by two key bottlenecks. First, existing datasets are not designed around fine-grained listener reactions. ViCo provides task-specific supervision but is limited in scale, while larger dyadic resources such as RealTalk~\cite{geng2023realtalk}, Seamless Interaction~\cite{agrawal2025seamless}, and SpeakerVid-5M~\cite{zhang2025speakervid5m} offer rich conversational videos without event-level annotations of reaction type, timing, and duration. As a result, current data provide insufficient supervision for modeling when a listener should react, what reaction should occur, and how the reaction should unfold over time.

Second, existing evaluation protocols mainly rely on metrics inherited from talking-head or video generation, such as reconstruction fidelity, perceptual realism, and motion quality. These metrics are useful for measuring global visual quality, but they do not explicitly assess whether generated listener reactions are type-consistent, temporally aligned with conversational cues, or visually realistic within reaction segments. This is particularly important because listener feedback is inherently multi-valid: an appropriate response may not exactly match a single reference timestamp, yet its type, timing, and quality should still be evaluated in a reaction-aware manner~\cite{de2012survey}.

To address these limitations, we revisit listening head generation from the perspectives of data, modeling, and evaluation. We curate a listening-head-specific dataset from RealTalk and Seamless Interaction by extracting aligned speaker-listener portrait pairs and annotating listener reactions as categorized temporal events. The resulting dataset contains approximately 147 hours of paired video data and 64,557 reaction instances. Based on this dataset, we introduce an audio-driven listening head generation method GLARE with prosody conditioning and a temporal reaction loss, enabling the model to better capture speaker-side cues and generate appropriate listener reactions at the right moment. We further propose reaction-centric evaluation metrics that measure reaction occurrence, temporal alignment, asymmetric timing deviation, and reaction-region visual quality.

Our main contributions are as follows:

\begin{enumerate}[leftmargin=1.5em, itemsep=0em, topsep=0em]
    \item We curate a listening-head dataset built from dyadic conversational videos, with aligned video-audio pairs and event-level annotations of listener reactions, comprising approximately 147 hours of video data and 64,557 reaction instances.
    \item We introduce an audio-driven reacting listener, namely GLARE, for listening head generation, equipped with a dedicated prosody conditioning mechanism and a temporal reaction loss.
    \item We propose reaction-oriented evaluation metrics for listening head generation that jointly measures visual fidelity, reaction accuracy, and temporal alignment, enabling finer-grained assessment of conversational appropriateness with extensive experiments comparing existing methods.
\end{enumerate}